\section*{Chapitre \ref{ch:b3:nuclear-physics} --- Physique nucléaire}

\begin{solution}{exo:b3:nuclear-physics:1}
(a) $^{12}$C~: $6, 6, 12$~; $^{14}$C~: $6, 8, 14$~; $^{235}$U~:
$92, 143, 235$~; $^{238}$U~: $92, 146, 238$. (b) Les carbones~; les
uraniums. (c) La chimie, c'est le nuage électronique, fixé par $Z$~;
la stabilité, c'est l'équilibre $N$--$Z$ à l'intérieur. (d) Des
isobares --- reliés par une désintégration $\beta^-$~: $^{14}\text{C} \to {}^{14}\text{N}
+ \text{e}^- + \bar\nu$, le tic-tac de l'horloge au radiocarbone.
\end{solution}

\begin{solution}{exo:b3:nuclear-physics:2}
(a) $\rho = Am_{\text{u}}/\tfrac43\pi r_0^3A =
3m_{\text{u}}/4\pi r_0^3 \approx \qty{2.3e17}{kg/m^3}$~: $A$ se
simplifie. (b) $\qty{5}{mL} \to \qty{1.2e12}{kg}$ --- un milliard de
tonnes dans une cuillerée. (c) $(R/R_{\text{atome}})^3 \approx
(4.7\times10^{-5})^3 \approx 10^{-13}$~: l'atome est du vide avec un
grain lourd. (d) Un volume $\propto A$ signifie que chaque nucléon
réclame une place fixe et ne se lie qu'à ses voisins de contact~: la
force sature à la portée du \unit{fm}.
\end{solution}

\begin{solution}{exo:b3:nuclear-physics:3}
(a) $\Delta m = 2(1.007276 + 1.008665) - 4.001506 =
\qty{0.030376}{u}$~: $B = \qty{28.3}{MeV}$. (b) $B/A =
\qty{7.07}{MeV}$. (c) $28.3\,\text{MeV}/13.6\,\text{eV} \approx
2\times10^6$~: six ordres de grandeur --- d'où un combustible
nucléaire qui surclasse le chimique d'un million. (d) L'$\alpha$ est
un colis préassemblé et exceptionnellement bon marché~: l'émettre
exporte quatre nucléons contre une ristourne de \qty{28.3}{MeV},
qu'aucun nucléon isolé ne peut égaler.
\end{solution}

\begin{solution}{exo:b3:nuclear-physics:4}
(a) $N = N_{\text{A}}/226 = \qty{2.66e21}{}$. (b) $\lambda =
\ln2/(1600\times\qty{3.156e7}{s}) = \qty{1.37e-11}{s^{-1}}$. (c)
$\mathcal A = \lambda N = \qty{3.7e10}{Bq}$~: un curie, par
construction --- le gramme des Curie eux-mêmes. (d) Trois demi-vies~:
$\qty{0.125}{g}$ de radium (le reste étant désormais du radon et ses
descendants, en marche vers le plomb).
\end{solution}

\begin{solution}{exo:b3:nuclear-physics:5}
(a) Volume $+882$~; surface $-261$~; Coulomb $-0.711\times676/3.83
= -126$~; asymétrie $-23.7\times16/56 = -6.8$ (\unit{MeV}). (b)
Avec l'appariement $+12/\sqrt{56} = +1.6$~: $B \approx \qty{490}{MeV}$,
$B/A = 8.76$ contre les $8.79$ mesurés --- trois pour mille d'écart,
à partir d'une goutte de liquide. (c) Le volume contre
surface$+$Coulomb~: le croisement de leurs vitesses de croissance est
ce qui sculpte le maximum au fer. (d) Coulomb est de longue portée~:
chaque proton repousse tous les autres, $\propto Z^2$, tandis que la
force forte, saturante, ne paie jamais que $\propto A$ --- la défaite
finale des poids lourds.
\end{solution}

\begin{solution}{exo:b3:nuclear-physics:6}
(a) La dérivée $\partial/\partial Z$ des termes de Coulomb et
d'asymétrie donne
$2a_{\text{C}}Z/A^{1/3} = 4a_{\text{A}}(A - 2Z)/A$~; résoudre.
(b) $Z^* = 50.5/1.163 = 43.4$~: le choix de la nature à $A = 101$ est
le ruthénium, $Z = 44$. (c) Sans Coulomb, $Z^* = A/2$~: deux échelles
de Fermi (\cref{ch:b3:quantum-statistics}), au meilleur prix
remplies à hauteur égale. (d) La place du fragment dans la vallée est
$Z^* \approx 41$~: cinq protons de moins --- une cascade de cinq
désintégrations $\beta^-$ l'y ramène, chacune convertissant un
neutron et émettant un électron et un antineutrino.
\end{solution}

\begin{solution}{exo:b3:nuclear-physics:7}
(a) $t = T_{1/2}\ln(1/0.3)/\ln2 = 5730\times1.74 \approx
\qty{e4}{}$ ans~: un feu de fin de période glaciaire. (b) $N_{^{14}\text{C}}
= (N_{\text{A}}/12)\times10^{-12} = \qty{5.0e10}{}$~; $\lambda =
\qty{3.8e-12}{s^{-1}}$~: $\mathcal A \approx \qty{0.2}{Bq}$ ---
une douzaine de désintégrations par minute et par gramme, le taux de
comptage de travail de tout laboratoire de datation. (c) Après neuf
demi-vies l'activité passe sous le bruit de fond~: l'horloge tourne
encore mais ne se lit plus. (d) Les roches n'ont jamais respiré le
carbone de l'atmosphère~; leurs horloges sont les primordiales ---
uranium--plomb, potassium--argon --- aux demi-vies de milliards
d'années.
\end{solution}

\begin{solution}{exo:b3:nuclear-physics:8}
(a) Les \qty{8.78}{MeV} de l'$\alpha$ sont bien sous le rebord de
\qty{26}{MeV}~: classiquement il ne peut ni sortir ni être entré ---
et pourtant la demi-vie vaut \qty{0.3}{\micro s}. (b)
$\lambda_1 = \ln2/\qty{3e-7}{s} = \qty{2.3e6}{s^{-1}}$~;
$\lambda_2 = \ln2/\qty{1.4e17}{s} = \qty{4.9e-18}{s^{-1}}$~:
rapport $\approx 5\times10^{23}$. (c) Le taux de passage tunnel vaut
$\eu^{-2G}$ avec l'exposant de Gamow $\propto Z/\sqrt E$~: diviser
$E$ par deux relève $2G$ d'une cinquantaine d'unités, et
$\eu^{55} \approx 10^{24}$ --- la scandaleuse droite de
Geiger--Nuttall, née d'une seule exponentielle. (d) Les protons
solaires passent aussi par effet tunnel~: la température du cœur doit
être juste assez haute pour que la queue de Maxwell multipliée par le
facteur de Gamow soutienne la combustion --- même falaise, gravie de
l'autre côté.
\end{solution}

\begin{solution}{exo:b3:nuclear-physics:9}
(a) La scission fait passer ${\sim}235$ nucléons d'une liaison de
\qty{7.6}{} à ${\sim}\qty{8.5}{MeV}$~: $235\times0.9 \approx
\qty{200}{MeV}$. (b) $N = \qty{2.56e24}{}$ noyaux~:
$E = \qty{8.2e13}{J} \approx 2000$ tonnes de pétrole --- par
kilogramme. (c) \qty{3}{GW} thermiques pendant un an font
\qty{9.5e16}{J}~: environ \qty{1.2}{t} d'$^{235}$U. (d) Les deux
fragments positifs jaillissent sous leur propre répulsion
coulombienne, emportant ${\sim}\qty{170}{MeV}$ d'énergie cinétique
qui devient de la chaleur en quelques micromètres --- un réacteur est
une bouilloire dont la flamme est un recul électrostatique.
\end{solution}

\begin{solution}{exo:b3:nuclear-physics:10}
(a) $10^4$ générations~: $(1.001)^{10^4} = \eu^{10} \approx
22000$ --- du mégawatt à la dizaine de gigawatts en une seconde~;
aucun moteur ne déplace une barre aussi vite. (b) Avec
$k - 1 < \beta$, la chaîne ne peut se reproduire sur les seuls
neutrons prompts~: chaque étape de croissance attend les traînards à
${\sim}\qty{10}{s}$, si bien que le temps de génération effectif est
d'une fraction de seconde ou plus. (c)
$(1.0005)^{10} \approx 1.005$~: un demi pour cent par seconde ---
le domaine de la barre et de l'opérateur. (d) Pousser $k - 1$
au-delà de $\beta$ a mis le réacteur sur l'horloge prompte à
\qty{e-4}{s}~: ils ont perdu le délai de grâce des neutrons retardés
qui rend les réacteurs pilotables.
\end{solution}

\begin{solution}{exo:b3:nuclear-physics:11}
(a) $\dot m = L_\odot/c^2 = \qty{4.2e9}{kg/s}$ de masse pure~; le
débit d'hydrogène $\dot m/0.0071 \approx \qty{6e11}{kg/s}$ ---
six cents millions de tonnes par seconde. (b) Hydrogène brûlable
${\sim}2\times10^{30}\times0.1\times0.75 = \qty{1.5e29}{kg}$~:
$t \approx \qty{2.5e17}{s} \approx 8$ milliards d'années --- le
Soleil est d'âge moyen. (c) Barrière ${\sim}\qty{1.4}{MeV}$ contre
$k_{\text{B}}T \approx \qty{1.3}{keV}$~: un déficit d'un facteur
mille --- seuls les protons les plus rapides de la queue de Maxwell,
par effet tunnel (\cref{exo:b3:nuclear-physics:8}), fusionnent~; d'où
un Soleil qui brûle pendant des milliards d'années au lieu
d'exploser. (d) La cendre est l'$^4$He --- stable, d'un contentement
doublement magique~; il n'y a pas de fragments riches en neutrons à
se désintégrer en $\beta$ pendant des siècles.
\end{solution}

\begin{solution}{exo:b3:nuclear-physics:12}
(a) $\lambda = \ln2/\qty{21600}{s} = \qty{3.2e-5}{s^{-1}}$~:
$N = \mathcal A/\lambda = \qty{1.6e13}{}$ noyaux ---
\qty{2.6}{ng}. De la médecine au nanogramme. (b) Six heures survivent
à l'examen mais non au week-end~: la dose s'éteint d'elle-même~; un
$\gamma$ pur de \qty{140}{keV} quitte le corps vers la caméra sans le
péage brûlant des $\alpha$ ou des $\beta$ pour les tissus. (c) Un
électron et un positon au repos s'annihilent en \emph{deux} photons
de \qty{511}{keV}, dos à dos (énergie et quantité de mouvement,
\cref{ch:b3:relativistic-dynamics})~: chaque paire en coïncidence
définit une droite passant par la tumeur, et des milliers de droites
la triangulent --- la tomographie par loi de conservation. (d)
\qty{2}{J/kg} réchauffent les tissus d'un demi-millikelvin ---
thermiquement, une gorgée d'eau tiède~; mais délivrés en
${\sim}10^{17}$ ionisations par kilogramme, chacune tranchant des
liaisons chimiques à l'échelle de l'\unit{eV}, c'est chimiquement
mortel pour une cellule qui se divise. Les joules mesurent la
chaleur~; les grays mesurent le sabotage.
\end{solution}

\begin{solution}{pb:b3:nuclear-physics:1}
\textbf{1.} $\qty{200}{MeV} = \qty{3.2e-11}{J}$~:
$20\times10^6/3.2\times10^{-11} = \qty{6.3e17}{}$ fissions par
seconde.
\textbf{2.} $5.4\times10^{22}$ fissions par jour $\times
235/N_{\text{A}}$~: environ \qty{21}{g} d'$^{235}$U par jour.
\textbf{3.} Charbon~: $\qty{1.7e12}{J}/\qty{3e7}{J/kg} \approx 58$
tonnes par jour --- un rapport de masses de $\sim3\times10^{6}$, qui
\emph{est} le rapport \unit{MeV} sur \unit{eV} des liaisons
nucléaires aux liaisons chimiques.
\textbf{4.} $235\times(8.5 - 7.6) \approx \qty{210}{MeV}$~:
cohérent, la petite différence étant portée au compte des neutrons et
des neutrinos.
\textbf{5.} Énergie cinétique des fragments $\to$ ionisation $\to$
chaleur, en quelques micromètres et microsecondes~; la conduction la
remet à l'eau --- le réacteur est une bouilloire qui arrête des
fragments.
\textbf{6.} Ils héritent du $N/Z \approx 1.55$ de l'uranium, bien
au-dessus du fond de la vallée à leur $A$~: chacun doit redescendre
par une chaîne de désintégrations $\beta^-$ --- une radioactivité
intégrée, par géométrie de la vallée.
\textbf{7.} $k$ = neutrons engendrés par neutron, une génération plus
tard~: $0.999$ s'éteint, $1.000$ tient (un réacteur en marche),
$1.001$ croît.
\textbf{8.} Les neutrons lents s'attardent près du noyau et l'appétit
de fission de l'$^{235}$U croît fortement aux énergies thermiques.
Les chocs élastiques cèdent l'énergie le plus vite à masses égales
--- et l'hydrogène égale la masse du neutron~: l'eau est un
modérateur fait sur mesure.
\textbf{9.} Le bore 10 dévore les neutrons ($\text{n} + {}^{10}
\text{B} \to \alpha + {}^{7}\text{Li}$)~: barres rentrées, neutrons
mangés, $k$ baisse~; barres sorties, $k$ monte --- la manette des
gaz.
\textbf{10.} $(1.0005)^{10^4} = \eu^{5} \approx 150$~: des puissances
multipliées par cent cinquante chaque seconde --- désespérant pour
de la mécanique.
\textbf{11.} Sur l'horloge retardée de \qty{0.1}{s}~:
$(1.0005)^{10} \approx 1.005$, un demi pour cent par seconde. Règle~:
garder $|k - 1|$ bien au-dessous de $\beta = 0.007$, pour que les
neutrons retardés aient toujours la voix prépondérante.
\textbf{12.} L'ébullition retire le modérateur~: les neutrons restent
rapides, la fission s'affame, $k$ baisse --- la conception modérée à
l'eau se bride elle-même (une contre-réaction négative qui manquait à
ses cousins au graphite).
\textbf{13.} La puissance résiduelle~: la radioactivité des fragments
--- la question 6 de la partie I --- obéit à des demi-vies, non à des
barres de contrôle, et donne encore des mégawatts juste après
l'arrêt.
\textbf{14.} Aller plus vite que la lumière dans l'eau~: $v > c/n$.
\textbf{15.} $v > 0.752c$~: $\gamma = 1.51$, donc $E_{\text{k}} >
0.51\times\qty{511}{keV} \approx \qty{0.26}{MeV}$.
\textbf{16.} Les électrons $\beta$ issus des désintégrations de
fragments portent des \unit{MeV}, et les $\gamma$ du cœur propulsent
par effet Compton des électrons à des énergies voisines~: les deux
dépassent largement \qty{0.26}{MeV} --- la piscine est pleine
d'électrons qualifiés.
\textbf{17.} Le spectre Tcherenkov se renforce vers les courtes
longueurs d'onde (une intensité $\propto1/\lambda^2$)~: l'œil reçoit
du bleu et du violet --- la couleur, c'est la pente du spectre.
\textbf{18.} L'eau atténue exponentiellement les $\gamma$ et les
neutrons~; huit mètres font des dizaines de longueurs de
demi-atténuation, si bien que la passerelle est au niveau du bruit de
fond. La lueur, c'est l'eau qui \emph{absorbe} l'énergie du
rayonnement --- la preuve visible que le blindage la mange.
\textbf{19.} Les fragments à vie courte meurent en quelques minutes
(l'affaiblissement)~; la traîne à vie plus longue --- la même
puissance résiduelle qu'à la question 13 --- entretient une faible
lueur et une vraie charge thermique pendant des jours.
\textbf{20.} Avec $T_{1/2} = \qty{66}{h}$, le parent perd un facteur
${\sim}6$ par semaine~: les stocks se désintègrent d'eux-mêmes, si
bien que les hôpitaux vivent d'une «~vache à technétium~»
hebdomadaire livrée par des réacteurs comme celui-ci.
\textbf{21.} $96/66 = 1.45$ demi-vie~: $2^{-1.45} \approx
0.36$ --- il reste environ \qty{36}{GBq}.
\textbf{22.} Thermalisés à température ambiante, $\lambda =
h/\sqrt{3m_{\text{n}}k_{\text{B}}T} \approx \qty{1.5}{\angstrom}$
(\cref{exo:b3:crystalline-solids:5})~: nés accordés aux espacements
réticulaires.
\textbf{23.} Cinq ans tuent toute demi-vie allant jusqu'au mois~:
l'activité et la puissance résiduelle chutent de plusieurs ordres de
grandeur, jusqu'à ce que l'élément puisse voyager dans un château de
plomb plutôt que dans une piscine.
\textbf{24.} ${\sim}\qty{200}{kW}$ --- une centaine de radiateurs
domestiques dans une piscine de cent tonnes~: quelques dizaines de
kelvins par jour au pire, évacués tranquillement par convection
naturelle --- une sûreté passive par pure capacité thermique.
\textbf{25.} $B/A$ paie \qty{200}{MeV} par scission et \qty{21}{g}
par jour~; les neutrons retardés prêtent les secondes qui rendent $k$
pilotable~; l'eau modère, refroidit, blinde --- et luit en bleu
précisément parce qu'elle travaille~; et le produit qui compte le
plus part en flacons pour les hôpitaux du lundi.
\end{solution}
